% Extracto editor-ready del contenido exclusivo de masas/16_hidrodinamica_y_gauge.tex.
% Las líneas 1--14 están publicadas byte-exactamente desde
% 74_rev2_realizaciones_i.tex:544--557. El propietario íntegro permanece
% preservado; este extracto conserva la continuación exclusiva.

\subsection{Comparison of the Electromagnetic and Chromodynamic Null Realizations}
The photon and the eight gluons belong to distinct null vectors. For the photon,

\[
m_\gamma=0
\]
in the electromagnetic branch. For the gluons,

\[
m_g=0
\]
in the adjoint representation of colour. Nullity is not obtained by taking the limit of a positive character; the domain bifurcates beforehand. Photon realization requires Maxwell theory, gauge invariance, and two transverse polarizations. Gluon realization requires the colour connection, its eight generators, and restriction to the physical gauge space.

Bidirectional propagation can be written

\[
c_\pm=\frac{c}{1\pm\epsilon},
\qquad
\epsilon=q_-^{90}-q_+^{90},
\]
with harmonic mean

\[
\frac{2}{1/c_++1/c_-}=c.
\]
The bidirectional velocity preserves \(c\); the unidirectional orientation preserves the bias. An additional gap in the genealogical record may be realized through Proca--Stueckelberg, but it must not be identified automatically with \(\epsilon\) or with the splitting of the action.

\Needspace{7\baselineskip}
